\setcounter{section}{1}

\Needspace{5\baselineskip}
\hypertarget{stochastic-gradient-langevin-dynamics}{%
\section{Stochastic Gradient Langevin Dynamics}\label{stochastic-gradient-langevin-dynamics}}

\Needspace{5\baselineskip}
\hypertarget{i.-mathematical-expression-of-stochastic-gradient-langevin-dynamics}{%
\subsection{I. Mathematical Expression of Stochastic Gradient Langevin Dynamics}\label{i.-mathematical-expression-of-stochastic-gradient-langevin-dynamics}}

Langevin dynamics originated in physics as a description of the random motion of particles in a potential energy field. In machine learning, stochastic gradient Langevin dynamics (SGLD) treats parameter optimization as the movement of a ``particle'' through an energy landscape (defined by the loss function). Its discretized update formula is as follows:

\Needspace{5\baselineskip}
Suppose we want to sample an image \(x\) from an unknown target real-data distribution \(p(x)\). Langevin dynamics provides an iterative formula:

\[
x_{i+1}=x_i+\frac{\gamma}{2}\nabla_x\log p(x_i)+\sqrt{\gamma}z_i
\]

\Needspace{5\baselineskip}
Here:

\begin{itemize}
\tightlist
\item
  \(x_i\) is the state at step \(i\) (an image vector).
\item
  \(\gamma\) is the step size.
\item
  \(z_i\sim\mathcal{N}(0,I)\) is standard Gaussian noise.
\item
  \(\nabla_x\log p(x_i)\) is called the score function: the gradient of the target distribution's log probability density function with respect to the input \(x\).
\end{itemize}

\Needspace{5\baselineskip}
This formula contains two competing forces:

\begin{enumerate}
\def\labelenumi{\arabic{enumi}.}
\tightlist
\item
  \textbf{Deterministic drift (gradient ascent)}: \(\frac{\gamma}{2}\nabla_x\log p(x_i)\) pulls the data point \(x_i\) in the direction of the fastest increase in probability density (uphill), making the image increasingly similar to a real image.
\item
  \textbf{Stochastic diffusion (noise injection)}: \(\sqrt{\gamma}z_i\) continually injects random perturbations. This not only prevents the model from becoming trapped in a local optimum (such as the ``hitting the tree'' mean point mentioned earlier), but also ensures that the final samples strictly follow the real \(p(x)\) distribution throughout the high-dimensional space, rather than collapsing to a single highest-probability point.
\end{enumerate}

\Needspace{5\baselineskip}
Although this may appear to be merely the addition of noise, it reflects a different objective from conventional gradient descent:

Gradient descent is an optimization algorithm whose objective is to find a (local or global) minimum of the loss function and obtain a point estimate, corresponding to maximum a posteriori (MAP) estimation.

SGLD is a Bayesian sampling algorithm. Its objective is not to find a ``best'' point, but to approximate the entire posterior distribution through sampling. This allows us to quantify uncertainty, that is, to understand the reliability of the parameter estimates.

In addition, with large datasets, computing the true gradient \(\nabla_x L(x_t)\) over the entire dataset is extremely expensive. SGLD approximates this gradient using a stochastic gradient \(\nabla_x L(x_t;\xi_t)\) computed from a mini-batch \(\xi_t\) to perform the calculations above. The resulting parameter sequence \(\{x_t\}\) also converges to the true posterior distribution \(p(x\mid\mathrm{data})\).

\Needspace{5\baselineskip}
Why this sampling procedure can approximate the posterior distribution:

\[
\Delta\theta_t=-\frac{\epsilon_t}{2}\nabla\tilde{U}(\theta_t)+\sqrt{\epsilon_t}\cdot\eta_t,\quad \eta_t\sim\mathcal{N}(0,I)
\]

\Needspace{5\baselineskip}
This formula contains two components:

\begin{enumerate}
\def\labelenumi{\arabic{enumi}.}
\tightlist
\item
  \(-\frac{\epsilon_t}{2}\nabla\tilde{U}(\theta_t)\): the stochastic gradient descent component. It attempts to push the parameters toward the point of highest posterior probability (the MAP estimate).
\item
  \(\sqrt{\epsilon_t}\cdot\eta_t\): the injected Gaussian noise component.
\end{enumerate}

\Needspace{5\baselineskip}
The crucial \(\sqrt{\epsilon}\) scaling:

Notice that the coefficient of the gradient is \(\epsilon\) (the step size), whereas the coefficient of the noise is \(\sqrt{\epsilon}\). As the step size \(\epsilon\to 0\), the noise term \(\sqrt{\epsilon}\) becomes much larger than the gradient term \(\epsilon\) (because \(0.0001\ll\sqrt{0.0001}=0.01\)).

This ensures that, even at the bottom of the valley (where the gradient is close to 0), noise still dominates and forces the parameters to keep moving.

\Needspace{5\baselineskip}
According to the derivation from the Fokker--Planck equation, the stationary distribution of this stochastic process converges exactly to:

\[
p(\theta)\propto \exp(-U(\theta))
\]

This is our posterior distribution.
In addition, with large datasets, computing the true gradient \(\nabla_x L(x_t)\) over the entire dataset is extremely expensive. SGLD approximates this gradient using a stochastic gradient \(\nabla_x L(x_t;\xi_t)\) computed from a mini-batch \(\xi_t\) to perform the calculations above. The resulting parameter sequence \(\{x_t\}\) also converges to the true posterior distribution \(p(x\mid\mathrm{data})\).

\Needspace{5\baselineskip}
\hypertarget{ii.-score-matching}{%
\subsection{II. Score Matching}\label{ii.-score-matching}}

\Needspace{5\baselineskip}
To run Langevin dynamics, we must know \(\nabla_x\log p(x)\). In practice, however, \(p(x)\) is extremely complex. A probability distribution can usually be written in energy-based form:

\[
p(x)=\frac{e^{-E_\theta(x)}}{Z_\theta}
\]

Here, \(Z_\theta=\int e^{-E_\theta(x)}dx\) is the partition function. In high-dimensional pixel space, summing or integrating over every possible image to compute \(Z_\theta\) is intractable.

\textbf{The useful property of the score function is that it eliminates \(Z_\theta\):}

\[
\nabla_x\log p(x)=\nabla_x(-E_\theta(x)-\log Z_\theta)=-\nabla_x E_\theta(x)
\]

\Needspace{5\baselineskip}
Since \(Z_\theta\) is a constant, its derivative with respect to \(x\) is simply 0. We can therefore train a neural network \(s_\theta(x)\) to fit this score directly:

\[
s_\theta(x)\approx\nabla_x\log p(x)
\]

This process is called score matching.

\Needspace{5\baselineskip}
\hypertarget{iii.-the-essence-of-diffusion-models}{%
\subsection{III. The Essence of Diffusion Models}\label{iii.-the-essence-of-diffusion-models}}

\Needspace{5\baselineskip}
In DDPM, the neural network \(\epsilon_\theta(x_t,t)\) predicts the noise \(\epsilon\). Mathematically, predicting noise and predicting the score are completely equivalent. The conversion between them is:

\[
\nabla_{x_t}\log p(x_t)=-\frac{\epsilon_\theta(x_t,t)}{\sqrt{1-\bar{\alpha}_t}}
\]

Thus, the denoising network in DDPM is essentially a score estimator. When DDPM uses the Euler--Maruyama method to solve the reverse SDE, each denoising operation it performs is essentially annealed Langevin dynamics with specific coefficients.

\Needspace{5\baselineskip}
The original DDPM reverse-sampling formula (generating \(x_{t-1}\) from \(x_t\)) is:

\[
x_{t-1}=\frac{1}{\sqrt{\alpha_t}}
\left(x_t-\frac{1-\alpha_t}{\sqrt{1-\bar{\alpha}_t}}\epsilon_\theta(x_t,t)\right)
+\sigma_t z
\]

\Needspace{5\baselineskip}
Now substitute the ``translation bridge'' above into this formula, replacing \(\epsilon_\theta\) with its gradient form. The formula becomes score-based:

\[
x_{t-1}=\frac{1}{\sqrt{\alpha_t}}x_t
+\frac{1-\alpha_t}{\sqrt{\alpha_t}}\nabla_{x_t}\log p_t(x_t)
+\sigma_t z
\]

\Needspace{5\baselineskip}
For comparison, the standard Langevin dynamics iteration is:

\[
x_{i+1}=x_i+\frac{\gamma}{2}\nabla_x\log p(x_i)+\sqrt{\gamma}\epsilon
\]

\Needspace{5\baselineskip}
The term-by-term correspondence between the two is now clear:

\Needspace{4\baselineskip}
\begin{enumerate}
\def\labelenumi{\arabic{enumi}.}
\tightlist
\item
  \textbf{Gradient-guidance term (deterministic drift)}
\end{enumerate}

Langevin dynamics: \(\frac{\gamma}{2}\nabla_x\log p(x_i)\). This represents ascent along the probability-density gradient (toward the data manifold), with step size \(\frac{\gamma}{2}\).

Diffusion model: \(\frac{1-\alpha_t}{\sqrt{\alpha_t}}\nabla_{x_t}\log p_t(x_t)\). This is precisely the force with which DDPM uses the neural network to push the image toward the real distribution at each denoising step. Here, \(\frac{1-\alpha_t}{\sqrt{\alpha_t}}\) is exactly equivalent to the Langevin dynamics step-size coefficient \(\frac{\gamma}{2}\). As \(t\) decreases, this step size also decays dynamically (annealing).

\Needspace{4\baselineskip}
\begin{enumerate}
\def\labelenumi{\arabic{enumi}.}
\setcounter{enumi}{1}
\tightlist
\item
  \textbf{Noise-injection term (random walk)}
\end{enumerate}

Langevin dynamics: \(\sqrt{\gamma}\epsilon\). This is a random perturbation that prevents trapping in a local optimum and preserves distributional properties.

Diffusion model: \(\sigma_t z\), where \(z\sim\mathcal{N}(0,I)\). In DDPM, the variance \(\sigma_t^2\) is usually set to \(\beta_t\) or \(\tilde{\beta}_t\). This corresponds exactly to the random-injection component of Langevin sampling. The physical force that breaks the ``multimodal mean (hitting the tree)'' deadlock in embodied AI comes from this term.

\Needspace{4\baselineskip}
\begin{enumerate}
\def\labelenumi{\arabic{enumi}.}
\setcounter{enumi}{2}
\tightlist
\item
  \textbf{State-scaling term (damping/momentum)}
\end{enumerate}

Langevin dynamics: \(x_i\) (with coefficient 1).

Diffusion model: \(\frac{1}{\sqrt{\alpha_t}}x_t\).

\textbf{Why is there a difference?} Because the DDPM forward noise-addition process is not simply the addition of noise; it is an Ornstein--Uhlenbeck (OU) process. During the forward process, DDPM scales the image down at each step (multiplying it by \(\sqrt{\alpha_t}\)) to prevent pixel variance from exploding over time. During reverse Langevin sampling, it must therefore divide by \(\sqrt{\alpha_t}\) to compensate for that scaling. This is equivalent to Langevin dynamics with ``damping'' or ``centripetal force'' in physics.

\Needspace{5\baselineskip}
\hypertarget{iv.-differential-equation-representation}{%
\subsection{IV. Differential-Equation Representation}\label{iv.-differential-equation-representation}}

Traditional DDPM (such as the formulation proposed by Ho et al.) uses discrete steps. When the number of noise-addition and denoising steps \(T\to\infty\), the discrete Markov chain becomes a continuous-time stochastic differential equation (SDE).

\Needspace{4\baselineskip}
\begin{enumerate}
\def\labelenumi{\arabic{enumi}.}
\tightlist
\item
  \textbf{Forward noise-addition process (corrupting the data)}
\end{enumerate}

\Needspace{5\baselineskip}
The forward process can be described by the following SDE:

\[
dx=f(x,t)dt+g(t)dw
\]

Here, \(w\) is standard Brownian motion (a continuous-time random walk), \(f(x,t)\) is the drift coefficient, and \(g(t)\) is the diffusion coefficient. This process requires no network; it turns a clear image \(x_0\) into pure noise \(x_T\) solely through physical rules.

\Needspace{4\baselineskip}
\begin{enumerate}
\def\labelenumi{\arabic{enumi}.}
\setcounter{enumi}{1}
\tightlist
\item
  \textbf{Reverse generation process (the continuous manifestation of Langevin dynamics)}
\end{enumerate}

\Needspace{5\baselineskip}
According to Anderson's theorem, any forward SDE of the form above has a corresponding reverse SDE that generates real data \(x_0\) backward from noise \(x_T\):

\[
dx=\left[f(x,t)-g(t)^2\nabla_x\log p_t(x)\right]dt+g(t)d\bar{w}
\]

Here, \(d\bar{w}\) is Brownian motion under time reversal. Consider the core term of this reverse SDE: \(-g(t)^2\nabla_x\log p_t(x)\). It corresponds exactly to the gradient-guidance term in Langevin dynamics.

\FloatBarrier
\Needspace{5\baselineskip}
\hypertarget{references-129}{%
\subsection{References}\label{references-129}}

\vspace{0.25em}
\begin{itemize}
\tightlist
\item
  Welling, M., \& Teh, Y. W. (2011). \href{https://www.stats.ox.ac.uk/~teh/research/compstats/WelTeh2011a.pdf}{Bayesian Learning via Stochastic Gradient Langevin Dynamics}. ICML.
\item
  Roberts, G. O., \& Tweedie, R. L. (1996). \href{https://doi.org/10.2307/3318418}{Exponential Convergence of Langevin Distributions and Their Discrete Approximations}. Bernoulli.
\item
  Hyvärinen, A. (2005). \href{https://www.jmlr.org/papers/v6/hyvarinen05a.html}{Estimation of Non-Normalized Statistical Models by Score Matching}. Journal of Machine Learning Research, 6, 695--709.
\item
  Ho, J., Jain, A., \& Abbeel, P. (2020). \href{https://arxiv.org/abs/2006.11239}{Denoising Diffusion Probabilistic Models}. NeurIPS.
\item
  Song, Y., Sohl-Dickstein, J., Kingma, D. P., Kumar, A., Ermon, S., \& Poole, B. (2021). \href{https://arxiv.org/abs/2011.13456}{Score-Based Generative Modeling through Stochastic Differential Equations}. ICLR.
\end{itemize}

\clearpage
